% ৪.২ Assembly Language
\section{Assembly Language}

\begin{learningbox}
এই টপিক শেষে শিক্ষার্থী assembly language, mnemonic code ও assembler-এর ভূমিকা ব্যাখ্যা করতে পারবে।
\end{learningbox}

\begin{definitionbox}{Assembly language}
যে low-level language-এ binary instruction-এর বদলে mnemonic code যেমন \texttt{ADD}, \texttt{MOV}, \texttt{SUB} ব্যবহার করা হয়, তাকে assembly language বলা হয়।
\end{definitionbox}

\begin{center}
\fbox{\parbox{0.88\textwidth}{
\centering
\textbf{Translation placeholder}\\[4pt]
Assembly source $\rightarrow$ Assembler $\rightarrow$ Machine code/Object code
}}
\end{center}

\begin{center}
\begin{tabular}{|p{0.26\textwidth}|p{0.30\textwidth}|p{0.30\textwidth}|}
\hline
\textbf{ভিত্তি} & \textbf{Machine} & \textbf{Assembly} \\
\hline
রূপ & 0 ও 1 & mnemonic code \\
\hline
Translator & লাগে না & assembler লাগে \\
\hline
বোঝা & খুব কঠিন & তুলনামূলক সহজ \\
\hline
\end{tabular}
\end{center}
